\documentclass[10pt,a4paper]{article}

\usepackage[utf8]{inputenc}
\usepackage[margin=0.75in]{geometry}
\usepackage{amsmath,amssymb}
\usepackage{booktabs}
\usepackage{tabularx}
\usepackage{multirow}
\usepackage{makecell}
\usepackage{caption}
\usepackage{microtype}
\usepackage{xcolor}
\usepackage{hyperref}

\hypersetup{
    colorlinks=true,
    linkcolor=blue!70!black,
    citecolor=blue!70!black,
    urlcolor=blue!70!black
}

\captionsetup[table]{
    font=small,
    labelfont=bf,
    skip=4pt,
    justification=centering
}

\title{\textbf{\Large Multi-Horizon Spatiotemporal Air Quality Forecasting: Empirical Benchmark Report}\\[0.4ex]
\large Performance Evaluation of Deep Learning, Machine Learning, and Climate Baselines (2019--2024)}
\author{\textbf{Adarsh Pradeep} \\ Department of Computer Science \& Engineering}
\date{September 2026}

\begin{document}

\maketitle

\begin{abstract}
This report provides the rigorous empirical benchmark of our multi-horizon air quality forecasting architecture evaluated across the state of Tamil Nadu, India. Using multimodal earth observation data from Sentinel-2 MSI and Sentinel-5P TROPOMI, models are evaluated on an independent holdout period (2023--2024, comprising $>4\times 10^6$ grid-cell observations). We compare deep spatiotemporal neural networks (ConvGRU Decoupled ST-ResUNet), regularized tabular models, and climate teleconnection baselines across two operational forecast horizons: short-term ($1$--$5$ days, $H=1$) and long-term ($150$ days, $H=30$).
\end{abstract}

\vspace{0.8em}
\section{Short-Term Forecasting Benchmark ($1$--$5$ Days, $H=1$)}
At short lead times, the objective is tracking rapid fluid advection, diurnal urban emissions, and localized plume formation. Table~\ref{tab:short_term} summarizes the test performance across $\text{NO}_2$, $\text{CO}$, and $\text{SO}_2$.

\begin{table}[htbp]
\centering
\small
\caption{Test Results for Short-Term ($1$--$5$ Days Ahead, $H=1$) on 2023--2024 Holdout Period.}
\label{tab:short_term}
\begin{tabularx}{\textwidth}{llccccccX}
\toprule
\textbf{Pollutant} & \textbf{Model Architecture} & \textbf{Methodology} & \textbf{$R^2$} & \textbf{RMSE ($\text{mol/m}^2$)} & \textbf{Pearson $r$} & \textbf{SSIM} & \textbf{CSI ($Q_{90}$)} & \textbf{Benchmark Remarks} \\
\midrule
\multirow{4}{*}{\textbf{$\text{NO}_2$}} 
& Persistence Baseline ($t_0$) & Autoregressive & 0.1852 & $1.037 \times 10^{-5}$ & 0.6299 & 0.6439 & 0.2984 & Rapid autocorrelation decay after 24h \\
& Monthly Spatial Climatology & Empirical Climatology & 0.5316 & $7.862 \times 10^{-6}$ & 0.7403 & 0.7356 & 0.2987 & Strong seasonal baseline floor \\
& \textbf{ConvGRU Decoupled ST-ResUNet} & \textbf{Spatiotemporal DL} & \textbf{0.6510} & $8.054 \times 10^{-6}$ & \textbf{0.8100} & \textbf{0.8964} & \textbf{0.3614} & \textbf{Optimal: Superior edge fidelity \& plume tracking} \\
& LightGBM (Spatial Gradients) & Gradient Boosted Trees & 0.5479 & $\mathbf{7.724 \times 10^{-6}}$ & 0.7486 & 0.7293 & 0.3190 & Tabular benchmark; lacks spatial convolutional priors \\
\midrule
\multirow{4}{*}{\textbf{$\text{CO}$}} 
& Persistence Baseline ($t_0$) & Autoregressive & $-0.6431$ & $8.702 \times 10^{-3}$ & 0.5548 & 0.5960 & 0.1298 & Negative skill ($R^2 < 0$) across synoptic shifts \\
& Monthly Spatial Climatology & Empirical Climatology & 0.6830 & $3.823 \times 10^{-3}$ & 0.8344 & 0.7117 & 0.0941 & Captures stationary regional background \\
& ConvGRU Decoupled ST-ResUNet & Spatiotemporal DL & 0.6213 & $4.129 \times 10^{-3}$ & 0.7991 & 0.6075 & 0.1627 & High correlation ($r \approx 0.80$) with slight spatial smoothing \\
& \textbf{Multi-Scale Regularized Ridge} & \textbf{Regularized Linear} & \textbf{0.6927} & $\mathbf{3.764 \times 10^{-3}}$ & \textbf{0.8352} & \textbf{0.7190} & 0.0907 & \textbf{Optimal: Highest overall $R^2$ and lowest mass bias (0.81\%)} \\
\midrule
\multirow{4}{*}{\textbf{$\text{SO}_2$}} 
& Persistence Baseline ($t_0$) & Autoregressive & $-0.6987$ & $1.770 \times 10^{-4}$ & 0.1565 & 0.2272 & 0.0817 & Uncorrelated due to intermittent point emissions \\
& Monthly Spatial Climatology & Empirical Climatology & 0.0592 & $\mathbf{1.317 \times 10^{-4}}$ & 0.2845 & 0.2886 & 0.0841 & Background mean only \\
& ConvGRU Decoupled ST-ResUNet & Spatiotemporal DL & 0.0715 & $1.362 \times 10^{-4}$ & \textbf{0.3146} & \textbf{0.5761} & 0.0871 & Best continuous spatial field mapping \\
& \textbf{Two-Stage Hurdle Classifier} & \textbf{Risk / Exceedance} & \textbf{N/A} & $1.378 \times 10^{-4}$ & 0.2750 & 0.2556 & \textbf{0.1109} & \textbf{Optimal: Plume detection (ROC-AUC = 0.6156, Recall = 53.5\%)} \\
\bottomrule
\end{tabularx}
\end{table}

\vspace{0.8em}
\section{Long-Term Sub-Seasonal Benchmark ($150$ Days, $H=30$)}
At $150$ days ahead, meteorological predictability is depleted. Deep recurrent neural networks suffer severe degradation due to parameter over-capacity on $6$ years of satellite observations ($5$--$6$ annual cycles). Table~\ref{tab:long_term} presents the comparative evaluation.

\begin{table}[htbp]
\centering
\small
\caption{Test Results for Long-Term ($150$ Days Ahead, $H=30$) on 2023--2024 Holdout Period.}
\label{tab:long_term}
\begin{tabularx}{\textwidth}{llccccccX}
\toprule
\textbf{Pollutant} & \textbf{Model Architecture} & \textbf{Methodology} & \textbf{$R^2$} & \textbf{RMSE ($\text{mol/m}^2$)} & \textbf{Pearson $r$} & \textbf{SSIM} & \textbf{CSI / PR-AUC} & \textbf{Benchmark Remarks} \\
\midrule
\multirow{4}{*}{\textbf{$\text{NO}_2$}} 
& Persistence Baseline ($t_0$) & Autoregressive & $-0.0180$ & $1.159 \times 10^{-5}$ & 0.5115 & 0.6072 & 0.2261 & Negative predictive skill; zero weather memory \\
& ConvGRU (Deep Learning) & Recurrent Network & 0.2104 & $1.021 \times 10^{-5}$ & 0.4851 & 0.6412 & 0.1824 & Degraded by parameter over-capacity \\
& Regularized Ridge Regression & Regularized Linear & 0.5170 & $7.983 \times 10^{-6}$ & 0.7393 & 0.7341 & 0.2903 & Approaches climatological limit \\
& \textbf{Monthly Spatial Climatology} & \textbf{Empirical Climatology} & \textbf{0.5316} & $\mathbf{7.862 \times 10^{-6}}$ & \textbf{0.7403} & \textbf{0.7356} & \textbf{0.2987} & \textbf{Optimal: Stationary, drift-free empirical ceiling} \\
\midrule
\multirow{4}{*}{\textbf{$\text{CO}$}} 
& Persistence Baseline ($t_0$) & Autoregressive & $-0.5238$ & $8.380 \times 10^{-3}$ & 0.5224 & 0.5776 & 0.0966 & Negative predictive skill \\
& ConvGRU (Deep Learning) & Recurrent Network & $\mathbf{-0.3412}$ & $7.612 \times 10^{-3}$ & 0.3415 & 0.5821 & 0.0512 & \textbf{Severe degradation ($R^2 < 0$) via recursive error accumulation} \\
& Monthly Spatial Climatology & Empirical Climatology & 0.6830 & $3.823 \times 10^{-3}$ & 0.8344 & 0.7117 & 0.0941 & Consistent seasonal baseline \\
& \textbf{Harmonic + ENSO/IOD Modes} & \textbf{Teleconnection ML} & \textbf{0.6872} & $\mathbf{3.797 \times 10^{-3}}$ & \textbf{0.8350} & \textbf{0.7106} & \textbf{0.0861} & \textbf{Optimal: Incorporates large-scale climate indices (ENSO/IOD)} \\
\midrule
\multirow{4}{*}{\textbf{$\text{SO}_2$}} 
& Persistence Baseline ($t_0$) & Autoregressive & $-0.6722$ & $1.756 \times 10^{-4}$ & 0.0962 & 0.2289 & 0.0709 & Zero skill; loss of stationarity \\
& ConvGRU (Deep Learning) & Recurrent Network & $\mathbf{-0.5824}$ & $1.890 \times 10^{-4}$ & 0.1241 & 0.4820 & 0.0410 & \textbf{Severe regression failure on highly sparse point emissions} \\
& Monthly Spatial Climatology & Empirical Climatology & 0.0592 & $\mathbf{1.317 \times 10^{-4}}$ & 0.2845 & 0.2886 & 0.0841 & Background mean baseline \\
& \textbf{HistGradientBoosting Classifier} & \textbf{Cost-Sensitive Risk} & \textbf{N/A} & $1.370 \times 10^{-4}$ & \textbf{0.2896} & \textbf{0.3041} & \textbf{PR-AUC = 0.3974} & \textbf{Optimal: Actionable 5-month plume outlook (ROC-AUC = 0.6174)} \\
\bottomrule
\end{tabularx}
\end{table}

\vspace{0.8em}
\section{Recommended Operational Architecture Across Horizons}
Table~\ref{tab:champion} synthesizes the optimal deployed model per pollutant and lead time, establishing a robust multi-horizon operational framework.

\begin{table}[htbp]
\centering
\small
\caption{Recommended Operational Architecture by Forecast Horizon.}
\label{tab:champion}
\begin{tabularx}{\textwidth}{llllX}
\toprule
\textbf{Horizon} & \textbf{Pollutant} & \textbf{Recommended Architecture} & \textbf{Model Category} & \textbf{Key Operational Justification} \\
\midrule
\multirow{3}{*}{\textbf{$1$--$5$ Days}} 
& $\text{NO}_2$ & ConvGRU Decoupled ST-ResUNet & Spatiotemporal DL & Preserves spatial gradient structures ($\text{SSIM} = 0.8964$) and plume transport. \\
& $\text{CO}$ & Multi-Scale Regularized Ridge & Regularized Linear & Outperforms recurrent networks ($R^2 = 0.6927$ vs $0.6213$) with minimal mass bias (0.81\%). \\
& $\text{SO}_2$ & Two-Stage Hurdle Classifier & Risk / Exceedance & Addresses extreme physical sparsity ($\sim 98\%$ zeros) without regression bias. \\
\midrule
\multirow{3}{*}{\textbf{$150$ Days}} 
& $\text{NO}_2$ & Monthly Spatial Climatology & Empirical Climatology & Atmospheric memory is zero at 150 days; empirical climatology provides the true ceiling. \\
& $\text{CO}$ & Harmonic + ENSO/IOD Regression & Teleconnection ML & Prevents recurrent network breakdown by coupling to macro oceanic-atmospheric modes. \\
& $\text{SO}_2$ & HistGradientBoosting Classifier & Cost-Sensitive Risk & Reliable early warnings of plume outbreaks ($\text{CSI} = 0.2337$) where regression fails. \\
\bottomrule
\end{tabularx}
\end{table}

\section{Discussion \& Defense Points}
\begin{enumerate}
    \item \textbf{Atmospheric Predictability Limit ($7$--$10$ Days):} Numerical and recurrent deep learning models depend on deterministic initial conditions. Beyond $10$ days, chaotic turbulence eliminates initial state memory, causing unconstrained recurrent architectures to drift.
    \item \textbf{Degrees of Freedom vs. Sample Size:} Sentinel-5P has only $6$ years of continuous data ($2019$--$2024$), corresponding to merely $5$--$6$ annual seasonal realizations. High-capacity neural networks (millions of parameters) overfit on this regime, while low-parameter harmonic and climatological formulations generalize robustly.
    \item \textbf{Decoupled Hazard Modeling for $\text{SO}_2$:} Treating $\text{SO}_2$ as continuous regression yields near-zero predictions due to extreme physical sparsity ($>98\%$ background). Reframing long-range prediction as extreme hazard classification yields operationally actionable early warnings.
\end{enumerate}

\end{document}
